\subsection{Numerical Methods and Verification of 1D Arterial Solvers}
\label{subsec:verification-literature}

The numerical-method state of the art for 1D arterial flow is shaped by the
hyperbolic balance-law structure of the area--flow equations. Finite-volume
methods remain attractive because they track conservation and discontinuity-like
features through numerical fluxes, while benchmark studies have also compared
finite element, discontinuous Galerkin, finite difference, and simplified
network-oriented schemes
\cite{LeVeque2002FiniteVolumeMethodsForHyperbolicProblems,BoileauEtAl2015Benchmark,Wang2014BloodFlowNetworks}.
In recent open implementations, MUSCL finite-volume discretizations and
Lax--Friedrichs-type fluxes are common because they give robust low-cost
baselines with simple boundary and coupling treatments
\cite{BenemeritoEtAl2024OpenBF,BeckersKolbe2025LaxFriedrichsHemodynamics}.

High-order and equilibrium-aware methods address different limitations of such
baselines. DG and WENO reconstructions can improve smooth-region resolution, but
their evidence depends on source quadrature, boundary treatment, positivity
handling, and whether limiter fallback or projections alter the formal order.
Well-balanced schemes target a separate requirement: a physically meaningful
steady or rest state should remain steady under the discrete operator. Recent
work on well-balanced explicit, semi-implicit, and high-order blood-flow
schemes makes this a central numerical property for elastic vessels and network
applications
\cite{LuccaEtAl2025FFRPrediction,PimentelGarciaEtAl2025HighOrderWellBalanced}.
That literature is directly relevant to this report because the selected
Rusanov/MUSCL realization is deliberately reported as not exactly well balanced
for the stenotic geometry-rest state.

Verification asks whether the implementation solves the declared equations and
produces reproducible numerical records. For hyperbolic finite-volume solvers,
relevant evidence includes conservation diagnostics, CFL reporting,
positivity/limiting records, grid and timestep refinement, boundary-regime
checks, exact equilibria, and manufactured solutions
\cite{LeVeque2002FiniteVolumeMethodsForHyperbolicProblems}. Arterial solver
benchmarks provide a useful cross-code reference for wave propagation and
network behavior, but they do not replace problem-specific verification of a
selected stenosis realization
\cite{BoileauEtAl2015Benchmark,BenemeritoEtAl2024OpenBF}.

This distinction matters for the present report. A benchmark-stage completion
row, an output-consistency check, or a cross-integrator discrepancy can be
useful evidence about the computational realization, but it is weaker than a
manufactured-solution convergence table or a demonstrated discrete equilibrium.
A fixed-area characteristic boundary rule also requires evidence that the sign
condition under which it is applied holds in the sampled run. Verification
therefore has to persist diagnostic quantities, not only final plots.

The literature also separates method order from study validity. A high-order
scheme can still be the wrong evidence for a pressure-ratio claim if the
pressure observation maps, wall model, inlet/outlet histories, and comparison
data are not matched. Conversely, a lower-order method can still support a
bounded implementation claim if its equations, numerical record, diagnostics,
and limitations are explicit. The present report follows the second path: it
states the selected method, exposes its rest-state defect, and limits the
resolved comparison to velocity-output discrepancies.

\subsection{Cross-Dimensional Comparison and Validation}
\label{subsec:cross-dimensional-comparison-review}

A 1D--3D comparison is meaningful only after the compared observable has been
mapped to a common quantity. For velocity, this can mean reducing a resolved
3D velocity field to section averages, radial profiles, or point/slab samples.
For pressure and FFR-style quantities, the required maps are different and
must include proximal/distal pressure observations and time-averaging
conventions
\cite{Blanco2018FFR1D3D,LuccaEtAl2025FFRPrediction,ShaikhEtAl2025CTFFR}.
Agreement under one operator is not agreement under another.

Validation is still a separate evidential category. A comparison with another
simulation can identify cross-model discrepancies, but it does not establish
physical accuracy for the target question unless the reference model and its
boundary, wall, material, and history data are themselves accepted for that
purpose. The present work therefore uses cross-model velocity discrepancy
language and reserves \emph{error} for exact solutions, manufactured solutions,
or accepted reference data under matched conditions.
